ANALOG FILTER DESIGN - LECTURE 6

These formulas use analog rad/s $\Omega$ and $H_a(s)$, NOT discrete rad/sample $\omega$ and H(z).

PASSBAND AND STOPBAND SPECS

Passband tolerance $\delta_p$:

$$1-\delta_p\le |H_a(j\Omega)|\le1+\delta_p$$

Stopband tolerance $\delta_s$:

$$|H_a(j\Omega)|\le\delta_s$$

Passband/stopband attenuation in dB:

$$A_p=-20\log_{10}(1-\delta_p)$$

$$A_s=-20\log_{10}(\delta_s)$$

Passband edge $\Omega_p$; stopband edge $\Omega_s$. Transition ratio $k=\Omega_p/\Omega_s$.

BUTTERWORTH LOWPASS

Maximally flat passband.

$$|H_a(j\Omega)|^2=\frac{1}{1+(\Omega/\Omega_c)^{2N}}$$

At $\Omega=\Omega_c$:

$$|H_a|^2=\frac12$$

$$G(\Omega_c)=-3.0103\text{ dB}$$

The minimum order from dB specifications is:

$$N\ge\frac{\log_{10}\left(\frac{10^{A_s/10}-1}{10^{A_p/10}-1}\right)}{2\log_{10}(\Omega_s/\Omega_p)}$$

Round N UP to the next integer.

One order increase changes high-frequency attenuation by approximately 6 dB per octave.

CHEBYSHEV TYPE I

Ripple in passband; monotonic stopband.

$$|H_a(j\Omega)|^2=\frac{1}{1+\epsilon^2T_N^2(\Omega/\Omega_p)}$$

$$\epsilon=\sqrt{10^{A_p/10}-1}$$

For $|v|\le1$:

$$T_N(v)=\cos(N\arccos v)$$

For $v\ge1$:

$$T_N(v)=\cosh(N\cosh^{-1}v)$$

Calculator-friendly recursion:

$$T_0(v)=1,\quad T_1(v)=v$$

$$T_N(v)=2vT_{N-1}(v)-T_{N-2}(v)$$

Chebyshev-I order formula:

$$N\ge\frac{\cosh^{-1}\left(\sqrt{\frac{10^{A_s/10}-1}{10^{A_p/10}-1}}\right)}{\cosh^{-1}(\Omega_s/\Omega_p)}$$

Round UP.

CHEBYSHEV TYPE II

Monotonic passband; ripple in stopband. Unlike Type I, Type II can have finite-frequency transmission zeros.

ELLIPTIC FILTER

Ripple in both passband AND stopband. Sharp transition for a given order; poles and stopband zeros shape the response.

BESSEL FILTER

Designed for flat group delay / nearly linear phase, at the expense of slower magnitude rolloff.

$$H_a(s)=\frac{b_0}{B_N(s)}$$

Bessel polynomials (unnormalized):

$$B_0(s)=1,\quad B_1(s)=s+1$$

$$B_N(s)=(2N-1)B_{N-1}(s)+s^2B_{N-2}(s)$$

ANALOG FILTER FREQUENCY TRANSFORMS

Start from normalized lowpass prototype $H_{LP}(s)$.

Lowpass scaled to $\Omega_c$:

$$H(s)=H_{LP}(s/\Omega_c)$$

Highpass with edge $\Omega_c$:

$$H(s)=H_{LP}(\Omega_c/s)$$

Bandpass, with center $\Omega_0$ and bandwidth B:

$$H(s)=H_{LP}\left(\frac{s^2+\Omega_0^2}{Bs}\right)$$

$$\Omega_0=\sqrt{\Omega_1\Omega_2}$$

$$B=\Omega_2-\Omega_1$$

Bandstop:

$$H(s)=H_{LP}\left(\frac{Bs}{s^2+\Omega_0^2}\right)$$

Useful DSP conversion (bilinear transform):

$$s=\frac{2}{T_s}\frac{1-z^{-1}}{1+z^{-1}}$$

This maps analog poles from the left half-plane inside the digital unit circle. It warps frequency:

$$\Omega=\frac{2}{T_s}\tan(\omega/2)$$

A separate analog design may use Butterworth, Chebyshev, elliptic, or Bessel prototypes. Keep analog rad/s and digital rad/sample distinct.
